\subsection{Базові алгоритми просторового перетворення бібліотеки OpenCV}

Для повноцінної, об'єктивної та всебічної оцінки ефективності розробленого барицентричного алгоритму необхідно мати класичні методи для порівняння. З цією метою в програмний комплекс було інтегровано відкриту бібліотеку комп'ютерного зору OpenCV. Вона є визнаним індустріальним стандартом у галузі обробки цифрових зображень і містить глибоко оптимізовані реалізації традиційних алгоритмів.

У межах даного дипломного проєкту бібліотека OpenCV виконує дві ключові функції. По-перше, її базові модулі використовуються для виконання низькорівневих операцій вводу/виводу. Функції \texttt{cv::imread} та \texttt{cv::imwrite} забезпечують швидке декодування графічних файлів у матричний вигляд, зручний для попіксельної обробки у C++, а також подальше збереження результатів. Окрім цього, інструментарій OpenCV застосовується для управління кольоровими просторами.

По-друге, функціонал OpenCV (зокрема \texttt{cv::resize}) застосовується для отримання альтернативних результатів масштабування з метою їх подальшого порівняння. Оскільки ідеальні (еталонні) зображення для розрахунку метрик якості генеруються за допомогою середовища MATLAB, методи OpenCV виступають у ролі базових конкурентів. Для аналізу було обрано три алгоритми різного рівня складності:

\begin{itemize}
	\item \textbf{Інтерполяція за найближчим сусідом (\texttt{INTER\_NEAREST}):} Найпростіший та найшвидший алгоритм. Кожному пікселю цільового зображення присвоюється значення геометрично найближчого пікселя вихідного зображення. Алгоритм має мінімальну обчислювальну складність $O(1)$ на піксель, проте спричиняє сильні візуальні артефакти — «aliasing».
	
	\item \textbf{Білінійна інтерполяція (\texttt{INTER\_LINEAR}):} Базовий метод, що обчислює значення нового пікселя як зважену суму чотирьох найближчих сусідніх пікселів (околиця $2 \times 2$). Метод працює швидко і згладжує зубчасті краї, проте призводить до ефекту неприродного «blurring» на різких границях та втрати дрібних деталей.
	
	\item \textbf{Бікубічна інтерполяція (\texttt{INTER\_CUBIC}):} Складніший математичний апарат, який враховує околицю розміром $4 \times 4$ пікселі (16 сусідів) за допомогою кубічних сплайнів. Цей алгоритм зберігає значно більше деталей, але є більш ресурсомістким і схильним до виникнення «ringing artifact» — появи світлих або темних контурів навколо контрастних об'єктів.
	
\end{itemize}

Наявність цих різнопланових реалізацій дозволяє створити ідеальну експериментальну базу. Порівнюючи результати розробленої барицентричної інтерполяції та класичних методів OpenCV із еталонними зображеннями, можна максимально об'єктивно оцінити компроміс між швидкістю виконання та підсумковою якістю.